% ECONOMICS_PROBLEM: {"id": "E58-600", "title": "Communication and forward guidance", "jel_code": "E58", "source_id": "OP-0250"}
% !TeX program = xelatex
\documentclass[11pt,a4paper]{article}
\usepackage[margin=23mm]{geometry}
\usepackage{fontspec}
\setmainfont{Latin Modern Roman}
\usepackage{amsmath,amssymb,mathtools}
\usepackage{xcolor,enumitem,booktabs,longtable,array,xurl,hyperref}
\definecolor{muted}{HTML}{53636C}
\hypersetup{colorlinks=true,urlcolor=blue,linkcolor=blue}
\setlength{\parindent}{0pt}
\setlength{\parskip}{5pt}
\newcommand{\R}{\mathbb R}
\newcommand{\E}{\mathbb E}
\newcommand{\N}{\mathbb N}
\newcommand{\ind}{\mathbf 1}
\newcommand{\Var}{\operatorname{Var}}
\newcommand{\Cov}{\operatorname{Cov}}
\newcommand{\supp}{\operatorname{supp}}
\DeclareMathOperator*{\argmin}{arg\,min}
\DeclareMathOperator*{\argmax}{arg\,max}
\newcommand{\entrymeta}[3]{{\sffamily\small\color{muted}\textbf{\texttt{#1}}\quad|\quad #2\par #3\par}}
\newcommand{\Problemref}[1]{\texttt{#1}}
\newcommand{\JELref}[2]{\texttt{#2} (JEL \texttt{#1})}
\begin{document}
\section*{Communication and forward guidance}
\textbf{Problem ID:} \texttt{E58-600}\quad \textbf{JEL:} \texttt{E58}\par
% BEGIN ECONOMICS STATEMENT
\label{problem:OP-0250}
{\sffamily\small\color{muted}Original subject group: Inflation and monetary policy\par}
\label{definitions:problem:30007531}
\textbf{Audit classification:} Published research agenda. Forward-guidance design and evidence are institutional priorities; the micro experiment is editorial. Verification concerns the cited version, not first priority or proof that this exact editorial specification remains unsolved on October 8, 2026.
\subsection*{Definitions and precise research target}
\paragraph{Editorial mathematical specification.} Fix a population of households or firms and a central-bank communication date. Randomize message \(Z_i\) among a control, a policy-path announcement, and an announcement with a credibility cue whose content is predeclared. Let \(B_i(z)\) be the resulting vector of beliefs about rates, inflation, and economic activity, and \(C_i(z)\) actual subsequent consumption or investment measured administratively. Define the total communication effect
\[
\tau_C(z)=E[C_i(z)-C_i(0)]
\]
and belief response \(\tau_B(z)=E[B_i(z)-B_i(0)]\). Identify these effects by randomized assignment, accounting for attrition and spillovers. To estimate a causal belief-to-action channel using assignment as an instrument, specify a scalar mediator and the applicable first-stage, monotonicity, and exclusion restrictions, or use an explicit multivariate structural model: messages may convey central-bank information or alter confidence through pathways other than the measured beliefs.
Determine which message features produce persistent belief changes and actual behavior rather than immediate survey revisions. Compare effects across predetermined attention, liquidity, and trust measures, correcting for multiple subgroup comparisons. For extrapolation to aggregate forward guidance, specify equilibrium price feedback and interactions among treated agents; household-level treatment effects alone do not establish the aggregate response. The target is effective and credible communication in a declared environment, with explicit limits on transportability and mediation, rather than a universal forward-guidance coefficient.

Let assignment probabilities satisfy \(P(Z_i=z\mid h_i)>0\), and distinguish assigned message from actual exposure. Under random assignment, consistency, integrable outcomes, and no interference, the intention-to-treat contrast equals the corresponding observed mean difference. A general vector-valued belief response cannot be interpreted through a generic scalar monotonicity assumption. If a scalar binary belief mediator \(B_i\) is defined, exclusion, a nonzero first stage, and \(B_i(1)\ge B_i(0)\) identify a complier effect; continuous or multidimensional beliefs need separate structural restrictions or bounds. With aggregate announcements everyone observes, individual randomization only varies salience or delivery, not the macro policy path itself. A saturation design must specify exposure to others' treatment to identify spillovers.
\subsection*{Variants and assumption changes}
The following are \emph{editorial variants}, not additional published open conjectures.
\begin{enumerate}
\item \emph{Randomized delivery.} Keep public policy content identical and randomize salience or presentation. Identify the effect of receiving that content among compliers, separating it from the effect of changing the announced policy path.
\item \emph{Information versus commitment.} Randomize messages conveying identical policy-rate paths but different disclosures about the central bank's economic assessment. Estimate total effects, while using a predeclared structural model to distinguish information from commitment mechanisms.
\end{enumerate}
\subsection*{References and source audit}
\begin{itemize}
\item Janet L. Yellen. \emph{Macroeconomic Research After the Crisis}. 2016. Locator: Section Inflation Dynamics. Primary text checked. \url{https://www.federalreserve.gov/newsevents/speech/yellen20161014a.htm}.
\item European Central Bank. \emph{Monetary policy, strategy and implementation}. Undated. Locator: Forward guidance, selected areas of focus. Primary text checked. \url{https://www.ecb.europa.eu/pub/economic-research/research_agenda/monetary_policy/html/index.en.html}.
\end{itemize}
Forward-guidance design and evidence are institutional priorities; the micro experiment is editorial.
\begin{enumerate}\raggedright
\item\hypertarget{definitions:source:30007531:1}{} Janet L. Yellen. 2016. \emph{Macroeconomic Research After the Crisis}. Section Inflation Dynamics.
\url{https://www.federalreserve.gov/newsevents/speech/yellen20161014a.htm}.
\item\hypertarget{definitions:source:30007531:2}{} European Central Bank. Undated / date unverified. \emph{Monetary policy, strategy and implementation}. Forward guidance, selected areas of focus.
\url{https://www.ecb.europa.eu/pub/economic-research/research_agenda/monetary_policy/html/index.en.html}.
\end{enumerate}

\subsection*{Common conventions: Source mathematical conventions}

These conventions apply to each entry unless explicitly replaced.

\textbf{Models and identification.} A primitive model \(m\) specifies all probability laws, feasible choices, information, equilibrium and measurement rules used by its target. The admissible class \(\mathcal M\) is fixed independently of outcomes. If \(P_O(m)\) is its observable probability law and \(\tau(m)\) the target, its identified set at observable law \(P\) is
\[
 \mathcal I_\tau(P)=\{\tau(m):m\in\mathcal M,\ P_O(m)=P\}.
\]
Point identification means that this set is a singleton. An empty set signals incompatibility of data and assumptions. Coordinate bounds need not describe the joint set. Bounds are sharp only if every retained target is attainable or arbitrarily approachable by compatible models. Finite bounds require explicit restrictions on unobserved quantities; unrestricted confounding, hidden wealth, extrapolation or arbitrary equilibrium selection can yield uninformative sets.

\textbf{Causal interventions.} A potential outcome \(Y_i(p)\) is the outcome under a complete policy regime \(p\), including timing, financing, information and market spillovers. The target population and units are fixed before treatment. With interference, use the complete assignment vector or a justified exposure mapping; individual no-interference notation alone cannot identify an economy-wide intervention. A difference of mean potential outcomes does not require the cross-world joint distribution. Conditioning on post-treatment migration, survival, enrollment or employment changes the estimand and needs separate justification.

\textbf{Statistical statements.} An estimator, confidence set, forecast or test specifies a sampling law, model class and loss/coverage criterion. Uniform \((1-\alpha)\)-coverage means \(\inf_{m\in\mathcal M}\Pr_m\{\tau(m)\in C_n\}\geq1-\alpha\), or an explicitly asymptotic version. The whole parameter space trivially has coverage, so informativeness requires a stated width, risk or power objective. A forecasting improvement is relative to a named benchmark, information set and evaluation population.

\textbf{Equilibrium and optimization.} An equilibrium comprises feasible plans, optimality, beliefs where needed, budget constraints and all market/resource clearing. The primitives or a stated benchmark define each correspondence \(\mathcal E(p;m)\). Nonemptiness is assumed or proved, never inferred just from compact policy sets. A selected-equilibrium target uses a declared selection rule; a robust target quantifies over all permitted equilibria. Use suprema or infima unless attainment is established. An \(\varepsilon\)-optimal policy is feasible and within \(\varepsilon\) of the finite optimum value. Report an empty feasible set separately.

\textbf{Welfare and accounting.} Normative utility, interpersonal normalization, social weights, numeraire, discounting and the use of public receipts are primitives. Behavioral utility need not coincide with the welfare criterion. Household welfare includes ownership income and rents once; adding profits again to compensating variation generally double-counts them. A monetary equivalent is defined by an explicit transfer and price convention, with monotonicity and range conditions. Average effects, mechanism contributions, transfers and welfare losses are different targets. Infinite sums require convergence and dynamic feasible plans require terminal or no-Ponzi conditions.

\textbf{Source versions.} A working-paper date, revision date and journal publication year are distinguished. A DOI for a journal version is not automatically the DOI of an earlier manuscript. Locators refer to the stated version and printed pages where specified. A metadata-only check does not verify a claim about a full-text conclusion. Every entry's source subsection narrows the provenance claim accordingly.
% END ECONOMICS STATEMENT
\end{document}
